% !TEX root = ../main.tex

\section{从项目目录到 PDF}

在「文件」菜单中选择项目目录，可以在当前窗口打开，也可以新开一个窗口。文件树按目录展示源文件和图片，编译产生的中间文件会收在 \texttt{.build/} 中。拖动分隔线可以改变文件树、编辑器和预览栏的宽度。

选中主文件后，按 \texttt{Ctrl+S} 保存并编译。编译成功后，预览栏即时显示 PDF。有错误或警告时，点击问题列表中的条目可定位到对应源码。可在工具栏中选择 XeLaTeX、pdfLaTeX 或 LuaLaTeX。中文项目可以用 XeLaTeX 搭配 \texttt{ctex}。

\section{在源码与预览间定位}

双击源码，可在预览中定位该段内容的排版位置。图片作为浮动体移到别页时，双击插图代码可跳到图片所在位置。双击预览中的正文则会通过 SyncTeX 返回对应源码。这个双栏示例也可用来试验定位。

预览中的文字可直接选择。按住 \texttt{Ctrl} 滚动鼠标滚轮可缩放，按住 \texttt{Alt} 拖动可平移页面。点击文内链接会跳到目标位置，网页链接则在系统浏览器打开。

\newpage

\section{编辑多文件项目}

输入 \verb|\input{sec| 可按前缀补全项目路径。右击文件也能复制它相对于项目根目录的路径。文献引用会从已引入的 \texttt{.bib} 文件中查找键名，交叉引用会列出项目内已有的 \verb|\label|。

文件大纲可跳到章节。项目搜索支持正文和标签，提供精确匹配、通配符和整词选项。插入菜单可生成公式、图片、表格或引用的 LaTeX 结构。较长的源码行可用 \texttt{Alt+Z} 切换自动换行。

外部编辑器保存文件后，FastLaTeXer 会同步磁盘变更。若当前文件也有未保存的编辑，会提示选择载入磁盘版本或保留编辑内容。

\section{版本历史与导出}

Git 面板可浏览提交记录和改动，并为项目设置远程仓库。远程认证使用本机的 Git 凭据或 SSH 密钥。编译完成后，可以在预览栏直接导出 PDF，也可以通过「文件 → 准备提交」打包项目源码。

\section{试用这个示例}

打开示例目录（\texttt{samples/presentation}），选中 \texttt{main.tex}，再按 \texttt{Ctrl+S}。正文与配图放在独立文件中，可用来测试路径补全和文件切换。选择「设置 → 主题」可切换图\ref{fig:light}至图\ref{fig:furina-dark}所展示的六种配色。
